\section{Integral lower degrees after the three-box exclusions}
\label{sec:alternating-four-tail-integral-cutoffs}

\paragraph{Status and scope.}
Root derivation HD-ALT-10. Complete root and separate
/root/sixfolds whole-proof reviews have passed. Its three-box
premises are now promoted after their own complete peer audits. This is an effective numerical
corollary of those exclusions and the already accepted shape bounds.
It is not counted as a separate strict reduction of the main gate.

In the non-Eisenstein ordinary transitive alternating-core problem,
transpose the label if necessary and write $\lambda=(n-k,\mu)$.
The all-count standard, two-box and three-box exclusions give $k\ge4$.
The accepted ALT07/ALT08 necessary inequalities are
\[
 k<\frac{1-\sqrt{3/5}}2n,\qquad
 \frac{k(n-1-k)}{(n-1)(n-2)}
       \le\frac1{5h}+\frac3{5e},\qquad
 e\ge\frac{n(n-3)}2.                                  \tag{ALT10.1}
\]
Here $h=2$ universally, $h=4$ if eight divides the original
quadratic discriminant, or $h=p-1$ for an odd discriminant
prime $p\ge5$. The actual inducing degree is $e$; equality
of the original field and the inducing character field is
not required by the cited shape inequalities.

\begin{corollary}[HD-ALT-10]
Every surviving label in this gate satisfies
\[
                              n\ge20h-2.                \tag{ALT10.2}
\]
This integer is optimal for the displayed numerical conditions:
at $n=20h-2$, the label $(n-4,4)$ satisfies them. This
numerical endpoint asserts no existence of a curve source.
\end{corollary}

\begin{proof}
If $n\le35$, the first inequality in (ALT10.1) implies
$k<35(1-\sqrt{3/5})/2<4$; the latter strict inequality
is equivalent to $3\cdot35^2>5\cdot27^2$.
Thus $n\ge36$. On $4\le k<n/4$, the expression
$k(n-1-k)$ is increasing, so it suffices to consider $k=4$.

Put $m=5h\ge10$ and suppose $36\le n\le4m-3$.
The strict contradiction to (ALT10.1) is
\[
 \frac{4(n-5)}{(n-1)(n-2)}-\frac1m
                              >\frac6{5n(n-3)}.         \tag{ALT10.3}
\]
After clearing positive denominators its numerator is
\[
 F(n,m)=5n(n-3)\bigl(4m(n-5)-(n-1)(n-2)\bigr)
                         -6m(n-1)(n-2).                 \tag{ALT10.4}
\]
Its coefficient of $m$ is positive for $n\ge36$:
$(n-3)(n-5)>n$ makes its first contribution exceed
$20n^2$, whereas the subtracted contribution is less
than $6n^2$. Because $m\ge(n+3)/4$, it follows that
\[
 \begin{split}
 F(n,m)&\ge F\bigl(n,(n+3)/4\bigr)\\
 &=\frac72n^3-100n^2+\frac{531}2n-9\\
 &=\frac72(n-36)^3+278(n-36)^2
                 +\frac{13347}2(n-36)+43245>0.
 \end{split}                                             \tag{ALT10.5}
\]
This proves (ALT10.3), which contradicts the degree floor
$3/(5e)\le6/(5n(n-3))$. Hence $n\ge4m-2=20h-2$.

At $n=4m-2$ the exact difference is
\[
 \frac{4(n-5)}{(n-1)(n-2)}-\frac1m
                    =\frac{-12}{m(n-1)(n-2)}<0.          \tag{ALT10.6}
\]
The known two-row degree formula, in the already cited
James--Liebeck Theorem~29.13 \cite{HD:JamesLiebeck2001}, gives
\[
 a=\binom n4-\binom n3=\frac{n(n-1)(n-2)(n-7)}{24},
 \qquad
 a-\frac{n(n-3)}2
       =\frac{n(n+1)(n^2-11n+22)}{24}>0                 \tag{ALT10.7}
\]
for $n\ge38$. The strict long-row condition also holds:
at $n=38$ it is $4<19(1-\sqrt{3/5})$, equivalent
to $3\cdot19^2<5\cdot15^2$, and its right side increases
with $n$. Thus the endpoint label has all the required
numerical properties, including for every positive tensor
power. Nothing in this verification supplies a generating
tuple, period-map rank or effective detecting cycle.
\end{proof}

For the eight non-Eisenstein class-number-one fields, the
resulting necessary lower bounds are
\[
\begin{array}{c|r|r}
K&h&n\text{ lower bound}\\\hline
\mathbf Q(i)&2&38\\
\mathbf Q(\sqrt{-2})&4&78\\
\mathbf Q(\sqrt{-7})&6&118\\
\mathbf Q(\sqrt{-11})&10&198\\
\mathbf Q(\sqrt{-19})&18&358\\
\mathbf Q(\sqrt{-43})&42&838\\
\mathbf Q(\sqrt{-67})&66&1318\\
\mathbf Q(\sqrt{-163})&162&3238
\end{array}                                               \tag{ALT10.8}
\]
This presents the reduced condition-A supplier gate
(cote16, scan050; SIGN\_NEXT SN-C03/C04) in effective integers,
with its existing finite-cover, local-system, representation
lambda-ring and K-theory interpretation. It changes no
cycle-supply premise and adds no Markman, Severn or032 input.
